\chapter{Manifiesto de LisPy}
\label{chap:manifesto}

LisPy nace de una convicción: no existe una razón técnica por la que un
programador deba elegir entre el poder expresivo de Lisp y el vasto
ecosistema de Python. La ingeniería de software moderna demanda ambos.

\section{La Visión}

LisPy integra dos runtimes en un mismo proceso: un intérprete de Common Lisp
escrito en C99 portable y el runtime embebido de CPython 3.14. El resultado
es un sistema donde las funciones de Lisp y Python conviven, se llaman
mutuamente y comparten memoria bajo un mismo techo de abstracción.

\section{¿Por qué Lisp?}

Lisp ofrece cualidades que ningún otro lenguaje iguala:

\begin{itemize}
\item \textbf{Expresividad máxima.} El código Lisp son datos (S-expressions),
  lo que permite metaprogramación sin precedentes mediante macros.
\item \textbf{Sistema de objetos (CLOS).} El Common Lisp Object System es el
  MOP (Meta-Object Protocol) más completo jamás diseñado, con métodos
  multimétodo, combinación de métodos y metamodelo reflejable.
\item \textbf{Sistema de condiciones.} El manejo de errores con
  \texttt{handler-case}, \texttt{handler-bind} y \texttt{restart-case}
  supera cualquier otro sistema de excepciones.
\item \textbf{Macros higiénicas.} La metaprogramación en tiempo de compilación
  permite crear DSLs (lenguajes de dominio específico) sin precedentes.
\end{itemize}

\section{¿Por qué Python?}

Python domina el panorama actual por razones igualmente sólidas:

\begin{itemize}
\item \textbf{Ecosistema.} NumPy, SciPy, scikit-learn, PyTorch, pandas,
  matplotlib -- ninguna plataforma científica iguala la madurez de Python.
\item \textbf{Comunidad.} Millones de desarrolladores, bibliotecas para todo.
\item \textbf{Rendimiento.} Aunque Python es lento como lenguaje, su ecosistema
  de extensiones nativas (C, Fortran, CUDA) ofrece velocidad nativa.
\end{itemize}

\section{¿Por qué Juntos?}

La combinación es cualitativamente superior a la suma de las partes. Lisp
orquesta y genera código; Python ejecuta cómputo intensivo. La comunicación
es bidireccional, con conversión automática de tipos y latencia
microscópica.

\section{Filosofía de Diseño}

LisPy se construye sobre principios claros:

\begin{itemize}
\item \textbf{Single-file C99 portable.} Todo el intérprete es un único
  archivo \texttt{lispy.c} de aproximadamente 66 000 líneas.
\item \textbf{CPython 3.14 embebido.} Sin DLLs externas ni dependencias
  frágiles: Python se linkea estáticamente.
\item \textbf{Puente bidireccional.} Lisp llama a Python y Python llama a
  Lisp con la misma naturalidad.
\item \textbf{Conversión automática.} Tipos Lisp y Python se convierten
  transparentemente siguiendo tablas bien definidas.
\end{itemize}

\section{Principio Rector}

\medskip
\begin{center}
\fbox{\parbox{0.8\textwidth}{\centering\itshape
Lisp como orquestador y generador,\\
Python como motor de cálculo.}}
\end{center}
\medskip

Este principio guía cada decisión de diseño: las partes del sistema que
requieren flexibilidad, metaprogramación y abstracción viven en Lisp. Las
partes que requieren cómputo numérico intensivo, acceso a bibliotecas
establecidas o procesamiento de datos delegan en Python.
